
\section{SA-ERM: Generic flat minima}\label{sec:sa_erm}
We begin by establishing a lower bound on the generalization performance of SA-ERM.
In particular, we construct an SCO instance where, with constant probability, there exists a minimizer of the SAER with a trivial $\Omega(1)$ population risk. This result illustrates not only the limitations of the SA-ERM algorithm in the general smooth SCO setting but also how the loss landscape affects generalization.
The result is formalized in the following theorem.
\begin{theorem}\label{thm: bad flat erm}
For every $n \in \mathbb{N}$ and $0 \leq \rho \leq \tfrac{1}{2}$, let $d = 2^n + 1$ and define 
$W = \{x \in \mathbb{R}^d : \|x\| \leq 1\}$. Then there exist an instance set $\mathcal{Z}$, a distribution $\mathcal{D}$ over $\mathcal{Z}$, and a loss function 
$f : W \times \mathcal{Z} \to \mathbb{R}$ that is convex, $1$-Lipschitz, $1$-smooth and $\rho$-flat, such that with probability at least $\tfrac{1}{2}$ over the training set $S$, there exist 
$w^{(1)}, w^{(2)} \in \arg\min_{w \in W} F_S(w)$ satisfying:
\begin{enumerate}[label=(\roman*),leftmargin=*,parsep=0pt,topsep=0pt]
    \item 
    for every $\rad \geq 0$, it holds that $w^{(1)} \in \arg\min_{w \in W} F_S^{\rad}(w)$. In particular, if $\rad \leq \rho$ then $w^{(1)}$ is an $\rad$-flat minimum of $F_S$; 
    \item 
    $w^{(2)}$ is a sharp minimum, in the sense that  $F_S^{\delta}(w^{(2)}) \geq F_S(w^{(2)}) + \frac{1}{2}\delta^2$ for all $\delta > 0$.%
    \footnote{This condition means that in every neighborhood of the minimizer there exists a point with large $F_S$. The inequality is the tightest possible: due to $1$-smoothness, any minimizer $w^\star$ of $F_S$ satisfies $F_S^{\delta}(w^\star) \leq F_S(w^\star) + \tfrac{1}{2}\delta^2$ for all $\delta > 0$.}
    \item 
    we have $F(w^{(1)}) - F(w^\star) = \Omega(1), \text{ while } F(w^{(2)}) - F(w^\star) = 0.$
\end{enumerate}
\end{theorem}

\cref{thm: bad flat erm} indicates that even when the loss is convex and $\beta$-smooth, and under the arguably strongest notion of flatness (\cref{ass:flat}), a flat minimum of the empirical risk may generalize poorly, whereas a sharp minimum of the same function can generalize optimally. We provide here a proof sketch, the full prove is deferred to \cref{sec:proof_sa_erm}.





\begin{proof}[sketch]
Our construction builds on classical lower bounds in stochastic convex optimization showing the existence of an empirical risk minimizers that overfits \citep{shalev2010learnability,feldman2016generalization}. In particular, \citet{shalev2010learnability} consider an instance space $Z = \{0,1\}^{d}$ with $d = 2^n$, where examples are drawn uniformly at random. Their loss function is of the form: 
\[
g(w,z) = \frac{1}{2} \sum_{i=1}^{d} z(i) w(i)^2 .
\]
With high probability, there exists a coordinate $I$ such that all sampled examples satisfy $z(I)=0$. The corresponding basis vector $e_I$ is then an ERM but incurs large population loss, yielding a spurious empirical minimizer.

A key difficulty in extending this construction is that the spurious ERM is not a flat minimizer, whereas SA-ERM favors flat solutions.
To address this, we use the observation that the function $h:\mathbb{R}\to\mathbb{R}$ defined by
\[
h(x) = \tfrac{1}{2}\max(x-\rho,0)^2
\]
is $1$-smooth and $\rho$-flat around its minimizers for $\rho \leq \tfrac{1}{2}$.
Composing this function with a suitable variant of the above construction yields a smooth SCO instance in which SA-ERM generalizes poorly. The resulting function is as follows:
\[
f(w,z) = \frac{1}{2}\left[ \sqrt{\sum_{i=1}^{d} z(i) w(i)^2 + w(d+1)^2} - \rho \right]_+^2 .
\]
\end{proof}




\section{SA-GD: Algorithmically chosen flat minima}\label{sec:sa_gd}
In the previous section, we presented in \cref{thm: bad flat erm} a hard instance showing that an existence of a flat minimizer that generalizes poorly. However, the same instance also admits flat minima that generalize well.
This raises a natural question: does this failure extend to practical algorithms explicitly designed to seek flat minima, such as SA-GD, or do such methods tend to converge to “good” flat minima? In this section we show that the former holds, and that the lower bound from \cref{thm: bad flat erm} also applies to flat minima obtained by SA-GD. 

We first establish a theorem on the optimization error of SA-GD, showing that when the perturbation radius $\rad$ is properly tuned, SA-GD minimizes the SAER objective and converges to a flat minimum of the empirical risk. 
The result is formalized in the following theorem whose proof is deferred to \cref{sec:proof_max_opt}.

\begin{theorem} \label{thm:max_opt}
Assume that $f(w,z)$ is $\beta$-smooth, convex, non-negative and $\rho$-flat for all $z$. Let $\{w_t\}_{t=1}^T$ be produced by SA-GD for $T$ steps
(\cref{gd_update_rule}) with $\eta \leq \ifrac{1}{4\beta}$ and $\rad >0$. For $\widehat{w} \coloneqq \frac{1}{T}\sum_{t=1}^T, w_t$ it
holds that
\begin{align*}
    F_S\left(\widehat{w}\right)\leq \fsmax\left(\widehat{w}\right)
    &\leq
    \frac{\|w_1-w^\star\|^2}{\eta T} + 4\beta\max\{\rad - \rho, 0\}^2
    .
\end{align*}
In particular, when $\eta=\ifrac{1}{4\beta}$, $\|w_1-w^\star\|=O(1)$ and $\rad-\rho=O(\ifrac{1}{\sqrt{T}})$, it holds that $$
F_S\left(\widehat{w}\right)\leq \fsmax\left(\widehat{w}\right)
    \leq
    O\left(\ifrac{\beta}{T}\right).$$
\end{theorem}

\cref{thm:max_opt} highlights the effect of flatness on the convergence rate of the algorithm.
When the flatness radius $\rho$ is small, the algorithm incurs an additive $O(\rad^2)$ term in the bound on the SAER objective. In contrast, when $\rho$ is large, even for $\rad \approx \rho$, SA-GD still minimizes the SAER objective and converges to a flat empirical minimum. Moreover, although SA-GD can be viewed as gradient descent applied to a potentially non-smooth function,\footnote{For example, if $F_S(x)=x^2$, which is $\beta$-smooth for $\beta=2$, then $\fsmax(x) = (|x|+\rad)^2$ is non-smooth.} its convergence rate in this case matches that of gradient descent on smooth functions.

For the proof, we first make use of the following key lemma, which establishes a regret bound for general algorithms whose update rule takes the form
$
w_{t+1} = w_t - \eta \nabla F_S(w_t + v_t)$, for $ \|v_t\| \leq \rad.
$

\begin{lemma}
\label{lem:gen_regret}
Assume that for every $z$, $f(w,z)$ is $\beta$-smooth, convex, non-negative and $\rho$-flat. Let $A$ be an algorithm that given a data set $S$, produces a sequence $\{w_t\}_{t=1}^T$ such that $
w_{t+1} = w_t - \eta \nabla F_S(w_t+v_t),$
 where $\{v_t\}_{t=1}^T$ are vectors such that for every $t$$, \|v_t\|\leq \rad$ and $\eta\leq 1/4\beta$. It holds that,
 \begin{equation*}
      \frac{1}{T}\sum_{i=1}^T F_S\left(w_t+v_t\right)-F_S(w^\star)  \leq 
    \frac{\|w_1-w^\star\|^2}{\eta T} + 4\beta\max\{\rad - \rho, 0\}^2.
 \end{equation*}
\end{lemma}

Next, we show that even when SA-GD converges to a flat minimum, the resulting solution is not guaranteed to generalize well. To demonstrate this, we establish the following lower bound on the population risk of SA-GD.
\begin{theorem}\label{thm: gen max sam}
    For every $n,T\in \mathbb{N},,\eta > 0, \rad \geq 0, \rho < \rad \bigl(1 - \frac{3}{3 + \eta \sqrt{T}}\bigr)$, assume $\eta (\rad - \rho) \leq \frac{1}{\sqrt{T}}$, let $d = 2^n T$ and define $W = \{x \in \mathbb{R}^d: \|x\| \leq 1\}$. Then there exists an instance set $\mathcal{Z}$, a distribution $\mathcal{D}$ over $\mathcal{Z}$, function $f: W \times \mathcal{Z} \to \mathbb{R}$ that is convex $1$-smooth, $1$-Lipschitz and $\rho$-flat, such that for a training set $S$ it holds that with probability at least $\frac{1}{2}$, running SA-GD for $T$ steps yields for every $\tau \in [T]$ suffix average $\widehat{w}_{\tau} = \frac{1}{T-\tau +1} \sum_{t=\tau}^T w_t$:
    $$F(\widehat{w}_{\tau}) - F(w^{\star}) = \Omega(\eta^2 (\rad - \rho)^2 T).$$
\end{theorem}

In particular, it follows that for step size $\eta \approx \ifrac{1}{\beta}$ and perturbation radius $\rad \approx \rho + \ifrac{1}{\sqrt{T}}$, the population risk of SA-GD can be as high as $\Omega(1)$, despite converging to a flat empirical minimum, as shown in \cref{thm:max_opt}.
This result extends the poor generalization result of flat minima given in \cref{thm: bad flat erm} also to SA-ERMs that is chosen algorithmically by a natural sharpness-aware gradient method. We provide here a proof sketch, the full proof is deferred to \cref{sec:proof_gen_max_sam}.


\begin{proof}[sketch]


The main technical challenge in the proof is that, in the non-smooth setting, prior constructions (e.g., \cite{amir2021sgd, koren2022benign, livni2024sample, schliserman2024dimension, vansover2025rapid}) exploit non-smoothness to shape the algorithm’s dynamics, whereas in the smooth setting such an approach is not possible. Instead, our key idea is to control the sequence of maximizers $\{v_t \in \arg\max_{\|v\| \leq \rad} F_S(w_t + v)\}_{t=1}^T$ to direct the dynamics toward a spurious ERM, and make sure that the sequence $\{v_t\}_{t=1}^T$ are aligned to such directions.
For this, we base our hard instance on the construction for SA-ERM given in \cref{thm: bad flat erm}. 
In that construction, in the first iteration we have $v_1 = \rad e_i$, where $e_i$ corresponds to the spurious ERM.
As a result, SA-GD makes a single step of size $\eta (\rad-\rho)$ toward this bad ERM. 
The remaining challenge is to ensure that the algorithm takes $T$ such steps in this direction. 
To achieve this, we replicate the construction across \(T\) mutually orthogonal subspaces, writing \(w=(w^{(1)},\dots,w^{(T)})\) with each block containing an independent copy of the hard instance. 
In this way, since $v_t$ is chosen in a different subspace at each iteration $t$, the algorithm makes a single step in each subspace and eventually converges to a bad ERM. The resulting loss function is:
\begin{align*}
& f(w,z) = \frac{1}{2}\left[  
\sqrt{
\sum_{i=1}^{2^n}\sum_{t=1}^T
[z(i)w^{(t)}(i)]_+^2
+ w(d)^2
}
-\rho
\right]_+^2.
\end{align*}
\end{proof}
Finally, we establish an upper bound for the population loss achieved by SA-GD in the following theorem.  

\begin{theorem}
\label{thm:max_gen_upper}
For every $z$, assume that $f(w,z)$ is
$\beta$-smooth, convex, non-negative and $\rho$-flat. Let $\{w_t\}_{t=1}^T$ be produced by SA-GD for $T$ steps
(\cref{gd_update_rule}) with $\eta \leq \ifrac{1}{4\beta}$ and $\rad >0$. For $\widehat{w} \coloneqq \frac{1}{T}\sum_{t=1}^T w_t$, it
holds that
\begin{align*}
    \mathbb E F (\widehat{w}) 
    \leq
    O& \bigg[\frac{\|w_1-w^\star\|^2}{\eta T} + \eta \beta^2 \rad  ^2  T+ \frac{\beta^2\eta T}{n^2} 
      +\bigg(\beta+\frac{\beta^3\eta^2 T^2}{n^2}\bigg)\max\{\rad - \rho, 0\}^2 \bigg]
    .
\end{align*}
In particular for  $T=n, \eta=O(\ifrac{1}{\beta})$, $\|w_1-w^\star\|=O(1)$ and $\rad-\rho=O(\ifrac{1}{\sqrt{T}})$ it holds that,
\begin{align*}
    \mathbb E F \left(\widehat{w}\right)= 
    O\left(\frac{\beta}{n}+ \beta\rad^2 n \right).
\end{align*}
\end{theorem}
We note that when $\rad = 0$, the bound in \cref{thm:max_gen_upper} coincides with the risk bounds of \citet{nikolakakis2022beyond, lei2020fine} for GD and SGD in convex, smooth, realizable settings. However, when $\rad > 0$, our bound contains an additional excess term of $\eta \beta^2 r^2 T$ compared to GD and SGD. This term nearly matches our lower bound in \cref{thm: gen max sam}, up to its dependence on $\eta$ and $\rho$.

The proof is deferred to \cref{sec:proof_max_gen_upper} and relies on a leave-one-out algorithmic stability argument. Specifically, we show that replacing a single training sample results in only a small change in the learned model, which in turn implies a small change in the loss. This stability property allows us to bound the generalization gap, that is, the difference between the empirical risk and the population loss \cite{bousquet2002stability,hardt2016train}. 



\section{SAM: Practically chosen flat minima} \label{sec:sam}

Finally, we analyze SAM, a well-studied and practically relevant algorithm introduced by \cite{foret2021sharpness} as a computationally efficient approximation of SA-GD. For this algorithm, we establish the following bound on the empirical risk. 
Similarly to SA-GD, the flatness of the empirical risk plays a significant role in the convergence of SAM, 
achieving fast convergence rates the function is $\rho$-flat and $\rad\leq \rho + \ifrac{1}{\sqrt{T}}$.

\begin{theorem}
\label{thm:asc_opt_upper}
For every $z$, assume that $f(w,z)$ is
$\beta$-smooth, convex, non-negative and $\rho$-flat. Let $\{w_t\}_{t=1}^T$ be produced by SAM for $T$ steps
(\cref{gd_update_rule}) with $\eta \leq \ifrac{1}{4\beta}$ and $\rad >0$. For $\widehat{w} \coloneqq \frac{1}{T}\sum_{t=1}^T w_t$, it
holds that
\begin{align*}
    F_S\left(\widehat{w}\right)
    &\leq
    \frac{\|w_1-w^\star\|^2}{\eta T} + 4\beta\max\{\rad - \rho, 0\}^2
    .
\end{align*}
In particular, if $\eta=\frac{1}{4\beta}$, $\|w_1-w^\star\|=O(1),\rad-\rho=O(\ifrac{1}{\sqrt{T}})$, $F_S\left(\widehat{w}\right)
    \leq
    O(\ifrac{\beta}{T}).$
\end{theorem}
The proof is deferred to \cref{proof_asc_opt_upper}. We note that \cref{thm:asc_opt_upper} establishes convergence rates in terms of the empirical risk. 
A natural question is whether SAM achieves similar rates
for the SAER. 
In the following theorem, we show that this is not the case: SAM might incur an additional term of $\Omega(\rad^2)$ in the convergence rate for the SAER, even for $\rho$-flat functions. 
This demonstrates that SAM can converge to a non-flat minimum, even when a $\rho$-flat minimum exists. 

\begin{theorem}\label{thm:ascent_opt_non_flate}
    For every $\eta> 0, n \in \mathbb{N}, \rad,\rho \leq \frac{1}{2}$, $W = [-1,1]$ there exists an instance set $\mathcal{Z}$ and a loss function $F_S:W \times \mathcal{Z} \to \mathbb{R}$ that is non-negative, convex, $1$-Lipschitz, $1$-smooth and $\rho$-flat such that running SAM on $F_S$ for $T$ steps holds, for any suffix average $\widehat{w}_{\tau} = \frac{1}{T-\tau+1}\sum_{t = \tau}^T w_t$,
    $$\forall ~ 0\leq\rad \leq \tfrac{1}{2}: \qquad \fsmax(\widehat{w}_{\tau}) - \fsmax(w^{\star}) = \Omega(\rad^2),$$
    that is, SAM converges to a sharp minimum.
\end{theorem}
The proof of \cref{thm:ascent_opt_non_flate} is deferred to \cref{sec:ascent_opt_non_flate}.
Finally, we turn to discuss the generalization guarantees of SAM.
In the following lower bound, we show that SAM can exhibit poor generalization in SCO under the realizable setting ($\rho = 0$), leaving the $\rho$-flat case ($\rho \gg 0$) for future work. 
\begin{theorem}\label{thm: gen asc sam}
    Given $n \geq 6 ,T \geq 6, \eta,\rad>0$ such that $\eta\rad \leq 1/2\sqrt{T}$, let $d = 2^n T$ and $W = \{w \in \mathbb{R}^d: \|w\| \leq 1\}$. 
    Then there exists an instance set $\mathcal{Z}$, a distribution $\mathcal{D}$ over $\mathcal{Z}$, a convex $6$-smooth $6$-Lipschitz and realizable function $f: W \times \mathcal{Z} \to \mathbb{R}$ such that for a training set $S$ with probability at least $\frac{1}{3}$ running SAM for $T$ steps with trajectory $\{w_t\}_{t=1}^T$, yields for every $\tau \in [T]$ suffix average $\widehat{w}_{\tau} = \frac{1}{T-\tau +1} \sum_{t=\tau}^T w_t$:
    \[F(\widehat{w}_{\tau}) - F(w^{\star}) = \Omega(\eta^2 \rad^2 T).\]
\end{theorem}
We provide here a proof sketch, the full proof is deferred to \cref{sec:proof_gen_asc_sam}.
\begin{proof}[sketch]
As with the construction for SA-GD in \cref{thm: gen max sam} we begin with a loss that admits a spurious ERM and $T$ orthogonal subspaces, 
$$f_1(w,z) = \frac12\sum_{i=1}^{2^n}\sum_{t=2}^{T} z(i)\,w^{(i)}(t)^2.$$ 
The main difficulty in this context is that the algorithm is initialized at $w_1 = 0$, which, in previous constructions,  is already a minimizer of the empirical risk. As a result, if we were to apply the same approach, SAM would remain at initialization throughout training and thus generalize well.  

To overcome this challenge, our key idea is to exploit the normalization of the ascent step, which can amplify small perturbations into meaningful progress. We begin by introducing a sufficiently small linear loss in the first orthogonal subspace,
\[
f_2(w,z) = \frac{\gamma}{2}\,[\,v_z^\top w^{(1)} + \delta_1\,]_+^{2},
\]
for an appropriate choice of \(v_z\) depending on the samples \(z\), with small \(\delta_1,\gamma>0\). 
Although the gradients of this function at the initialization point are small, the normalization step amplifies them, producing a progress of $\eta \rad$ toward the bad ERM in this subspace. To extend this effect across \(T\) orthogonal subspaces, we design a chaining mechanism that couples consecutive subspaces, such that progress in one subspace activates progress in the next. This is achieved via the following function:
\begin{align*}
    f_3(w) =  \tfrac{1}{2}\sum_{i=1}^{2^n}\sum_{t=2}^{T} 
\left[w^{(i)}(t)-\delta_t\cdot w^{(i)}(t-1),0\right]_+^{\!2},
\end{align*}
for an appropriate choice of \(\{\delta_t\}_{t=2}^T\). This construction ensures that the algorithm makes progress of order \(\eta \rad\) in each of \(T\) distinct subspaces, ultimately guiding the iterates toward a spurious ERM that generalizes poorly. 
The final loss function is therefore
\begin{align*}
f(w,z) = f_1(w,z) + f_2(w) + f_3(w,z). \tag*{\qedhere}
\end{align*}  
\end{proof}

Finally, we establish an upper bound for the population loss achieved by SAM in the following theorem. 

\begin{theorem}
\label{thm:ascent_gen_upper}
For every $z$, assume that $f(w,z)$ is
$\beta$-smooth, convex, non-negative and $\rho$-flat. Let $\{w_t\}_{t=1}^T$ be produced by SAM for $T$ steps
(\cref{asc_update_rule}) with $\eta \leq \ifrac{1}{4\beta}$ and $\rad >0$. For $\widehat{w} \coloneqq \frac{1}{T}\sum_{t=1}^T w_t$, it
holds that\begin{align*}
    \mathbb E F (\widehat{w}) 
    \leq
    O& \bigg[\frac{\|w_1-w^\star\|^2}{\eta T} + \eta \beta^2 \rad  ^2  T+ \frac{\beta^2\eta T}{n^2} +\bigg(\beta+\frac{\beta^3\eta^2 T^2}{n^2}\bigg)\max\{\rad - \rho, 0\}^2 \bigg]
    .
\end{align*}
In particular for  $T=n, \eta=O(\ifrac{1}{\beta})$, $\|w_1-w^\star\|=O(1)$ and $\rad-\rho=O(\ifrac{1}{\sqrt{T}})$ it holds that,
\begin{align*}
     \mathbb E F \left(\widehat{w}\right) =O\left(\frac{\beta}{n}+ \beta\rad^2 n \right).
\end{align*}
\end{theorem}

As in the bound for SA-GD, we note that when $\rad = 0$, the bound in \cref{thm:max_gen_upper} coincides with the risk bounds of \citet{ lei2020fine,nikolakakis2022beyond} for GD and SGD in convex, smooth, realizable settings. However, when $\rad > 0$, our bound contains an additional excess term of $\eta \beta^2 r^2 T$ relative to GD and SGD. This term nearly matches our lower bound in \cref{thm: gen asc sam}, up to its dependence on $\eta$.


The proof follows from an algorithmic stability analysis, similar to that in
\cref{thm:max_gen_upper};, and is given in \cref{sec:proof_ascent_gen_upper}.

